\section{Πρόβλημα 3}

Θεωρούμε ένα παίγνιο με άπειρο πλήθος αμιγών στρατηγικών, όπου ο ιδιοκτήτης μιας
εταιρείας (παίκτης 1) προσλαμβάνει έναν υπάλληλο (παίκτης 2) και του αναθέτει
κάποια δουλειά. Ο ιδιοκτήτης αποφασίζει για τον μισθό $x \in \mathbb{R}^+$ που
θα πάρει ο υπάλληλος, και ο υπάλληλος αποφασίζει για το επίπεδο προσπάθειας
$y \in \mathbb{R}^+$ που θα καταβάλει. Οι συναρτήσεις ωφέλειας των 2 παικτών
είναι:
$$
    u_1(x, y) = 2\sqrt{y} - (x + 2)^2 + 5xy
$$
$$
    u_2(x, y) = x - \frac{y^2}{2} + \alpha xy
$$
όπου $\alpha \geq 1$ είναι μια σταθερά.

Για να βρούμε τα σημεία ισορροπίας θα βρούμε τα μέγιστα των συναρτήσεων ωφέλειας
του κάθε παίκτη, θεωρώντας την άλλη μεταβλητή σταθερή. Για τον παίκτη 1 έχουμε
\begin{align*}
    & \frac{\partial u_1}{\partial x} = -2(x + 2) + 5y \\
\end{align*}
Παίρνωντας και την δεύτερη μερική παράγωγο, έχουμε
$$
    \frac{\partial^2 u_1}{\partial x^2} = -2 < 0 \,.
$$
Άρα για
\begin{align*}
    \frac{\partial u_1}{\partial x} &= 0 \\
    \Leftrightarrow x &= \frac{5y}{2} - 2
\end{align*}
η $u_1$ παίρνει την μέγιστη τιμή της.

Αντίστοιχα για τον παίκτη 2 έχουμε
\begin{align*}
    & \frac{\partial u_2}{\partial y} = - y + \alpha x \\
\end{align*}
Παίρνωντας και την δεύτερη μερική παράγωγο, έχουμε:
$$
    \frac{\partial^2 u_2}{\partial y^2} = -1 < 0 \,.
$$
Άρα για
\begin{align*}
    \frac{\partial u_2}{\partial y} &= 0 \\
    \Leftrightarrow y &= \alpha x
\end{align*}
η $u_2$ παίρνει την μέγιστη τιμή της.

Για να είναι σημείο ισορροπίας θέλουμε τα δύο μέγιστα να ικανοποιούνται
ταυτόχρονα, δηλαδή
$$
\begin{cases}
    x = \frac{5y}{2} - 2 \\
    y = \alpha x
\end{cases}
\Leftrightarrow
\begin{cases}
    x = \frac{5\alpha x}{2} - 2 \Leftrightarrow x(1-\frac{5 \alpha}{2}) = -2
    \Leftrightarrow x = \frac{4}{5\alpha - 2} \\
    y = \frac{4 \alpha}{5\alpha - 2}
\end{cases}
\,,
$$

Άρα για κάθε $\alpha \geq 1$ υπάρχει μοναδικό σημείο ισορροπίας το
$(x, y) = (\frac{4}{5\alpha - 2}, \frac{4 \alpha}{5\alpha - 2})$.